
\begin{subremarque}
\label{XIII.3.1.3}
Supongamos
\marginpar{409}
que $k=\kres(s)$ es de característica cero; entonces los esquemas de
tipo finito sobre~$\overline{k}$ de dimensión $\leqslant\dim X$ son
fuertemente desingularizables, y la demostración de~\Ref{XIII.3.1}
permite probar los resultados siguientes:
\begin{enumerate}
\item[a)] Existe un abierto no vacío~$S_1$ de~$S$ tal que, para todo
esquema~$S'$ sobre~$S_1$ cuyos puntos maximales son de característica
cero y para todo haz de conjuntos localmente constante y constructible
$F$ sobre~$X'=X\times_SS'$, $(F,f')$ sea cohomológicamente propio
relativamente a~$S'$ en dimensión $\leqslant 0$.

\item[b)] Si todas las características residuales de~$S$ son nulas,
existe un abierto no vacío~$S_1$ de~$S$ tal que, para todo esquema~$S'$
sobre~$S_1$ y para todo haz de grupos localmente constante y constructible
$F$ sobre~$X'$, $(F,f')$ sea cohomológicamente propio relativamente
a~$S'$ en dimensión $\leqslant 1$.
\end{enumerate}

En efecto, basta repetir la demostración de~\Ref{XIII.3.1}\ifthenelse{\boolean{orig}}
{.}
{~}%
2)\kern2pt a). La proposición~\Ref{XIII.2.7} utilizada en 3)\kern2pt 4
se aplica al caso de un haz localmente constante~$F$, pues, retomando las
notaciones de 3)\kern2pt4, todo torsor bajo~$F$ está moderadamente
ramificado sobre~$P$ relativamente a~$S$, ya que todas las características
residuales de~$S$ son nulas.
\end{subremarque}
